\documentclass[cjk]{loom}

\runningthread{online learning}

\begin{document}

\loomcover
  {The Adversary's Loom}
  {Notes on online classification \& the Littlestone dimension}
  {北极甜虾 (剑心犹在！)}
  {Michaelmas, 2026}

\section{The repeated game}

We play in rounds $t=1,\dots,T$. A \keyword{learner} must guess a hidden
labelling; an \keyword{adversary} reveals the truth one bite at a time and is
free to choose the worst continuation. Nothing is random. The only question is:
how many times can the learner be \emph{wrong}?

\begin{strand}
The picture to hold: the adversary is not a fixed distribution, it is a player
watching you. Every bound below is a statement about a \emph{tree of futures},
not about an average. If you remember one thing, remember the tree.
\end{strand}

\begin{definition}[mistake bound]
A learner $A$ for a class $\mathcal H \subseteq \{0,1\}^{\mathcal X}$ has
\keyword{mistake bound} $M$ if, against every realisable sequence
$(x_1,y_1),\dots,(x_T,y_T)$ with $y_t = h^\star(x_t)$ for some
$h^\star \in \mathcal H$, it errs at most $M$ times.
\end{definition}

\warp{ldim}%
\begin{definition}[Littlestone dimension]
$\operatorname{Ldim}(\mathcal H)$ is the depth of the deepest \emph{shattered}
binary tree: a complete tree whose internal nodes carry points of $\mathcal X$
and whose every root-to-leaf branch is realised by some $h \in \mathcal H$.
\end{definition}

\loose{Is the deepest tree unique? Almost never. Does the SOA need to find it?
No --- it only needs its \emph{depth}. Revisit.}%
The two definitions are secretly the same statement read from two sides. That
is the whole content of the next knot.

\section{The knot that ties it together}

\begin{theorem}[Littlestone, 1988]\label{thm:soa}
For every class $\mathcal H$, the optimal deterministic mistake bound equals
$\operatorname{Ldim}(\mathcal H)$. The bound is achieved by the
\keyword{Standard Optimal Algorithm} (SOA), which predicts the label keeping the
larger residual Littlestone dimension.
\end{theorem}

\begin{proof}
\emph{Upper bound.} Each time the SOA errs, the adversary's reply contradicts
the majority side, so the surviving subclass drops its Littlestone dimension by
at least one (this uses \pick{ldim}). Starting from $\operatorname{Ldim}(\mathcal
H)$, there can be at most that many mistakes.

\emph{Lower bound.} Walk the shattered tree: at each node force the SOA to
commit, then send the branch it did \emph{not} choose. Every step is a mistake,
and the tree is realisable by construction, so $\operatorname{Ldim}(\mathcal H)$
mistakes are unavoidable.
\end{proof}

\warmth{4} I can rebuild the upper bound cold; the lower bound's
``send the other branch'' still feels like a rabbit from a hat.

\begin{example}[thresholds on a line]
Let $\mathcal X = [0,1]$ and $\mathcal H$ the thresholds
$h_a(x)=\mathbf 1[x \ge a]$. A shattered tree of depth $k$ exists for every $k$
(binary search keeps halving the live interval), so
$\operatorname{Ldim} = \infty$. \emph{Online}, thresholds are unlearnable ---
yet their VC dimension is $1$ and they are trivially PAC-learnable.
\end{example}

\begin{strand}
中文直觉：在线与离线的鸿沟，全藏在那棵「未来之树」里。VC 维只数一层标注的自由度，
而 Littlestone 维数数的是\emph{对手能逼你走多深}。阈值函数恰好是这道裂缝最干净的见证。
\end{strand}

\begin{remark}
Randomisation softens but does not close the gap: the optimal expected mistake
bound sits between $\tfrac12\operatorname{Ldim}$ and $\operatorname{Ldim}$, and
the constant is itself a small open story.
\end{remark}

\loose{Open thread to pull next session: does a \emph{margin} assumption collapse
$\operatorname{Ldim}$ back down to something finite? Cf.\ the large-margin
perceptron bound $1/\gamma^2$.}%
The separation above is the seed of everything else in these notes: agnostic
regret, bandit feedback, and the apple-tasting variants all measure how badly
the tree can branch when the learner is handicapped.

\end{document}
